% ============================================================
%  RAPPORT DE FIN D'ETUDES — Chapitre 3
%  Module : Authentification et Gestion des Roles
%  Etudiant  : Ghoulam Mohamed Said
%  Equipe    : Reddas Belkacem · Chelig Mohamed Amine
%  Encadreur : Dr. Abdelkader Bougassa
%  Compatible Overleaf — pdfLaTeX
% ============================================================

\documentclass[12pt,a4paper,oneside]{report}

% ── Encodage & langue ──────────────────────────────────────
\usepackage[utf8]{inputenc}
\usepackage[T1]{fontenc}
\usepackage[french]{babel}
\usepackage{lmodern}

% ── Mise en page ───────────────────────────────────────────
\usepackage[top=2.5cm,bottom=2.5cm,left=3cm,right=2.5cm]{geometry}
\usepackage{setspace}
\onehalfspacing
\usepackage{parskip}
\setlength{\parindent}{0pt}

% ── Couleurs ───────────────────────────────────────────────
\usepackage{xcolor}
\definecolor{GoldenGlow}{RGB}{212,175,55}
\definecolor{SoftEspresso}{RGB}{62,39,35}
\definecolor{AccentGray}{RGB}{120,110,100}
\definecolor{CodeBg}{RGB}{248,246,240}

% ── Graphiques ─────────────────────────────────────────────
\usepackage{graphicx}
\usepackage{tikz}
\usetikzlibrary{shapes.geometric,arrows.meta,positioning,calc}
\usepackage{pgfplots}
\pgfplotsset{compat=1.18}

% ── Flottants & légendes ───────────────────────────────────
\usepackage{float}
\usepackage{caption}

% ── Tableaux ───────────────────────────────────────────────
\usepackage{booktabs}
\usepackage{tabularx}
\usepackage{array}
\newcolumntype{C}[1]{>{\centering\arraybackslash}p{#1}}
\newcolumntype{L}[1]{>{\raggedright\arraybackslash}p{#1}}

% ── Code source ────────────────────────────────────────────
\usepackage{listings}
\lstset{
  backgroundcolor=\color{CodeBg},
  basicstyle=\ttfamily\small,
  breaklines=true,
  captionpos=b,
  commentstyle=\color{AccentGray},
  frame=single,
  framerule=0.5pt,
  rulecolor=\color{GoldenGlow},
  keywordstyle=\color{SoftEspresso}\bfseries,
  stringstyle=\color{GoldenGlow!80!black},
  numbers=left,
  numberstyle=\tiny\color{AccentGray},
  numbersep=8pt,
  showspaces=false,
  showtabs=false,
  tabsize=2,
  literate=
    {é}{{\'{e}}}1 {è}{{\`{e}}}1 {ê}{{\^{e}}}1 {ë}{{\"e}}1
    {à}{{\`{a}}}1 {â}{{\^{a}}}1 {ù}{{\`{u}}}1 {û}{{\^{u}}}1
    {î}{{\^{\i}}}1 {ï}{{\"{\i}}}1 {ô}{{\^{o}}}1 {ç}{{\c{c}}}1
    {<}{{\textless}}1 {>}{{\textgreater}}1
    {->}{{$\rightarrow$}}2
}

% ── Boites colorées ────────────────────────────────────────
\usepackage[most]{tcolorbox}
\tcbuselibrary{breakable,skins}

\newtcolorbox{infobox}[1]{
  colback=GoldenGlow!8,
  colframe=GoldenGlow,
  fonttitle=\bfseries\color{SoftEspresso},
  title=#1,
  arc=4pt,
  boxrule=0.7pt,
  left=6pt,right=6pt,top=4pt,bottom=4pt
}

% ── Chapitres & sections ───────────────────────────────────
\usepackage{titlesec}
\titleformat{\chapter}[display]
  {\normalfont\huge\bfseries\color{SoftEspresso}}
  {\color{GoldenGlow}\chaptertitlename\ \thechapter}
  {10pt}{\Huge}
  [\vspace{4pt}\color{GoldenGlow}\titlerule[1.5pt]]

\titleformat{\section}
  {\normalfont\Large\bfseries\color{SoftEspresso}}
  {\color{GoldenGlow}\thesection}{1em}{}

\titleformat{\subsection}
  {\normalfont\large\bfseries\color{SoftEspresso!80}}
  {\thesubsection}{1em}{}

\titleformat{\subsubsection}
  {\normalfont\normalsize\bfseries\color{AccentGray}}
  {\thesubsubsection}{1em}{}

% ── En-tetes & pieds de page ───────────────────────────────
\usepackage{fancyhdr}
\pagestyle{fancy}
\fancyhf{}
\fancyhead[L]{\small\color{SoftEspresso}\nouppercase{\leftmark}}
\fancyhead[R]{\small\color{GoldenGlow}Authentification \& Gestion des Roles}
\fancyfoot[C]{\small\color{AccentGray}\thepage}
\renewcommand{\headrulewidth}{0.4pt}

% ── Utilitaires ────────────────────────────────────────────
\usepackage{enumitem}
\usepackage{amsmath}
\usepackage{microtype}
\usepackage{hyperref}
\hypersetup{
  colorlinks=true,
  linkcolor=SoftEspresso,
  urlcolor=GoldenGlow!80!black
}

% ── Commandes perso ────────────────────────────────────────
\newcommand{\role}[1]{\texttt{\color{GoldenGlow!80!black}#1}}
\newcommand{\endpoint}[1]{\texttt{\small\color{SoftEspresso}#1}}
\newcommand{\tech}[1]{\textbf{\color{SoftEspresso}#1}}

% ============================================================
\begin{document}
% ============================================================

\setcounter{chapter}{2}   % ce fichier est le Chapitre 3

% ════════════════════════════════════════════════════════════
\chapter{Implémentation et Tests}
% ════════════════════════════════════════════════════════════

\section{Introduction}

Ce chapitre décrit l'implémentation du module \textbf{Authentification et
Gestion des Rôles} du Système de Gestion Universitaire Intégré (SGUI). Nous
présentons successivement l'environnement de développement, les composants
backend (service JWT, middlewares, contrôleur), les composants frontend
(contexte d'authentification, routes protégées, page de connexion), puis les
résultats des tests fonctionnels et de sécurité.

% ────────────────────────────────────────────────────────────
\section{Environnement de Développement}
% ────────────────────────────────────────────────────────────

\begin{table}[H]
\centering
\caption{Environnement de développement — Module Authentification \& Rôles}
\begin{tabular}{@{}ll@{}}
\toprule
\textbf{Composant}        & \textbf{Technologie / Version}          \\
\midrule
Langage backend           & TypeScript 5.x                          \\
Runtime                   & Node.js 20 LTS                          \\
Framework API             & Express.js 4.18                         \\
ORM                       & Prisma 5.x                              \\
Base de données           & PostgreSQL 16                           \\
Hachage mot de passe      & bcryptjs 2.4.x (facteur de cout 12)     \\
Tokens JWT                & jsonwebtoken 9.x (algorithme HS256)     \\
Framework frontend        & React.js 18                             \\
Styles                    & Tailwind CSS 3.x                        \\
Client HTTP               & Axios                                   \\
\bottomrule
\end{tabular}
\end{table}

% ────────────────────────────────────────────────────────────
\section{Architecture du Module}
% ────────────────────────────────────────────────────────────

Le module suit une architecture en couches. Chaque requête traverse
obligatoirement les middlewares d'authentification et d'autorisation
avant d'atteindre le contrôleur métier.

\begin{figure}[H]
\centering
\begin{tikzpicture}[
  box/.style={
    draw=SoftEspresso, thick, rounded corners=5pt,
    minimum width=2.8cm, minimum height=0.9cm,
    fill=GoldenGlow!10, font=\small\bfseries, align=center
  },
  arrow/.style={->, >=Stealth, thick, SoftEspresso}
]
\node[box] (C)  at (0,   0) {Client\\React.js};
\node[box] (A)  at (3.8, 0) {requireAuth\\Middleware};
\node[box] (R)  at (7.6, 0) {requireRole\\Middleware};
\node[box] (Ct) at (11.4,0) {Auth\\Controller};
\node[box] (DB) at (7.6,-2) {PostgreSQL\\Prisma};

\draw[arrow] (C)  -- node[above,font=\scriptsize]{Bearer JWT} (A);
\draw[arrow] (A)  -- node[above,font=\scriptsize]{req.user}   (R);
\draw[arrow] (R)  -- node[above,font=\scriptsize]{next()}     (Ct);
\draw[arrow] (Ct) -- node[right,font=\scriptsize]{Prisma}     (DB);
\end{tikzpicture}
\caption{Pipeline de traitement d'une requête authentifiée}
\label{fig:auth_pipeline}
\end{figure}

\subsection{Organisation des Fichiers}

\begin{lstlisting}[language=bash,
  caption=Structure du module auth cote backend]
backend/src/
  modules/
    auth/
      auth.controller.ts    % Routes login / logout / me
      auth.service.ts       % Logique JWT + bcrypt
    users/
      users.controller.ts   % Gestion des comptes (admin)
  shared/
    middlewares/
      requireAuth.ts        % Verification du JWT
      requireRole.ts        % Verification du role
    prisma.ts               % Client Prisma singleton
\end{lstlisting}

% ────────────────────────────────────────────────────────────
\section{Modèle de Données}
% ────────────────────────────────────────────────────────────

\subsection{Schéma Prisma — Entité User}

\begin{lstlisting}[language=SQL,
  caption=schema.prisma --- Modele User et enumeration Role]
model User {
  id           String   @id @default(uuid())
  email        String   @unique
  passwordHash String
  role         Role
  nom          String
  prenom       String
  isActive     Boolean  @default(true)
  createdAt    DateTime @default(now())
  updatedAt    DateTime @updatedAt
}

enum Role {
  admin
  teacher
  student
}
\end{lstlisting}

\subsection{Matrice d'Autorisation}

\begin{table}[H]
\centering
\caption{Matrice d'autorisation des routes API}
\begin{tabularx}{\linewidth}{@{}llCCC@{}}
\toprule
\textbf{Methode} & \textbf{Route} &
  \textbf{admin} & \textbf{teacher} & \textbf{student} \\
\midrule
POST & \texttt{/api/auth/login}            & \checkmark & \checkmark & \checkmark \\
GET  & \texttt{/api/auth/me}               & \checkmark & \checkmark & \checkmark \\
POST & \texttt{/api/auth/logout}           & \checkmark & \checkmark & \checkmark \\
GET  & \texttt{/api/users}                 & \checkmark & $\times$   & $\times$   \\
POST & \texttt{/api/users}                 & \checkmark & $\times$   & $\times$   \\
GET  & \texttt{/api/pfe/sujets}            & \checkmark & \checkmark & \checkmark \\
POST & \texttt{/api/pfe/sujets}            & \checkmark & \checkmark & $\times$   \\
GET  & \texttt{/api/requests/admin/inbox}  & \checkmark & $\times$   & $\times$   \\
POST & \texttt{/api/requests/reclamations} & $\times$   & $\times$   & \checkmark \\
\bottomrule
\end{tabularx}
\end{table}

% ────────────────────────────────────────────────────────────
\section{Implémentation Backend}
% ────────────────────────────────────────────────────────────

\subsection{Service d'Authentification}

\begin{lstlisting}[language=JavaScript,
  caption=auth.service.ts --- Connexion et generation du token JWT]
import jwt    from 'jsonwebtoken';
import bcrypt from 'bcryptjs';
import { prisma } from '../shared/prisma';

const SECRET     = process.env.JWT_SECRET;
const EXPIRES_IN = '24h';

export async function loginUser(email, password) {
  const user = await prisma.user.findUnique({ where: { email } });

  if (!user || !user.isActive)
    throw new Error('INVALID_CREDENTIALS');

  const valid = await bcrypt.compare(password, user.passwordHash);
  if (!valid) throw new Error('INVALID_CREDENTIALS');

  const token = jwt.sign(
    { sub: user.id, role: user.role },
    SECRET,
    { expiresIn: EXPIRES_IN }
  );

  return {
    token,
    user: {
      id:     user.id,
      email:  user.email,
      nom:    user.nom,
      prenom: user.prenom,
      role:   user.role
    }
  };
}

export function verifyToken(token) {
  return jwt.verify(token, SECRET);
}
\end{lstlisting}

\subsection{Hachage des Mots de Passe}

Le hachage utilise \tech{bcrypt} avec un facteur de coût de~12, offrant
une résistance aux attaques par force brute tout en restant compatible
avec les performances d'un serveur standard (~250~ms par vérification).

\begin{lstlisting}[language=JavaScript,
  caption=auth.service.ts --- Creation d'un compte avec hachage bcrypt]
const SALT_ROUNDS = 12;

export async function createUser(email, password, role, nom, prenom) {
  const passwordHash = await bcrypt.hash(password, SALT_ROUNDS);

  return prisma.user.create({
    data: { email, passwordHash, nom, prenom,
            role, isActive: true }
  });
}
\end{lstlisting}

\subsection{Middleware d'Authentification}

\begin{lstlisting}[language=JavaScript,
  caption=middlewares/requireAuth.ts --- Verification du token JWT]
import { verifyToken } from '../modules/auth/auth.service';

export function requireAuth(req, res, next) {
  const header = req.headers.authorization;

  if (!header || !header.startsWith('Bearer '))
    return res.status(401).json({ error: 'TOKEN_MISSING' });

  try {
    const payload = verifyToken(header.split(' ')[1]);
    req.user = payload;   // { sub, role }
    next();
  } catch (err) {
    return res.status(401).json({ error: 'TOKEN_INVALID_OR_EXPIRED' });
  }
}
\end{lstlisting}

\subsection{Middleware d'Autorisation par Rôle}

\begin{lstlisting}[language=JavaScript,
  caption=middlewares/requireRole.ts --- Controle d'acces par role]
export function requireRole(roles) {
  return (req, res, next) => {
    const user = req.user;

    if (!user || !roles.includes(user.role))
      return res.status(403).json({ error: 'FORBIDDEN' });

    next();
  };
}
\end{lstlisting}

Exemple d'application sur les routes~:

\begin{lstlisting}[language=JavaScript,
  caption=Exemple d'application des middlewares sur les routes]
// Admin seulement
router.get('/users',
  requireAuth,
  requireRole(['admin']),
  listUsers
);

// Admin ET enseignant
router.post('/pfe/sujets',
  requireAuth,
  requireRole(['admin', 'teacher']),
  createSujet
);

// Tous les utilisateurs connectes
router.get('/auth/me',
  requireAuth,
  getProfile
);
\end{lstlisting}

\subsection{Contrôleur d'Authentification}

\begin{lstlisting}[language=JavaScript,
  caption=auth.controller.ts --- Routes login / profil / logout]
import { Router } from 'express';
import { loginUser } from './auth.service';
import { requireAuth } from '../../shared/middlewares/requireAuth';
import { prisma } from '../../shared/prisma';

const router = Router();

// POST /api/auth/login
router.post('/login', async (req, res) => {
  try {
    const { email, password } = req.body;
    if (!email || !password)
      return res.status(400).json({ error: 'MISSING_FIELDS' });

    const result = await loginUser(email, password);
    res.json(result);

  } catch (err) {
    if (err.message === 'INVALID_CREDENTIALS')
      return res.status(401).json({ error: 'INVALID_CREDENTIALS' });
    res.status(500).json({ error: 'INTERNAL_ERROR' });
  }
});

// GET /api/auth/me
router.get('/me', requireAuth, async (req, res) => {
  const user = await prisma.user.findUnique({
    where:  { id: req.user.sub },
    select: {
      id: true, email: true, role: true,
      nom: true, prenom: true, isActive: true
    }
  });
  if (!user) return res.status(404).json({ error: 'USER_NOT_FOUND' });
  res.json(user);
});

// POST /api/auth/logout
router.post('/logout', requireAuth, (req, res) => {
  res.json({ message: 'Deconnexion reussie.' });
});

export default router;
\end{lstlisting}

% ────────────────────────────────────────────────────────────
\section{Implémentation Frontend}
% ────────────────────────────────────────────────────────────

\subsection{Organisation des Fichiers}

\begin{lstlisting}[language=bash,
  caption=Structure frontend --- Module Authentification et Roles]
frontend/src/
  contexts/
    AuthContext.jsx      % Etat global d'authentification
  components/
    ProtectedRoute.jsx   % Garde de route par role
  services/
    api.js               % Instance Axios + intercepteurs
  pages/
    LoginPage.jsx        % Formulaire de connexion
\end{lstlisting}

\subsection{Contexte d'Authentification React}

Le \texttt{AuthContext} centralise l'état d'authentification et le rend
disponible à tous les composants via le hook \texttt{useAuth()}.

\begin{lstlisting}[language=JavaScript,
  caption=contexts/AuthContext.jsx --- Contexte global d'authentification]
import { createContext, useContext, useState, useEffect }
  from 'react';
import api from '../services/api';

const AuthContext = createContext(null);

export function AuthProvider({ children }) {
  const [user,    setUser]    = useState(null);
  const [loading, setLoading] = useState(true);

  useEffect(() => {
    const token = localStorage.getItem('token');
    if (token) {
      api.get('/auth/me')
        .then(r   => setUser(r.data))
        .catch(()  => localStorage.removeItem('token'))
        .finally(() => setLoading(false));
    } else {
      setLoading(false);
    }
  }, []);

  const login = async (email, password) => {
    const { data } = await api.post('/auth/login',
                                    { email, password });
    localStorage.setItem('token', data.token);
    setUser(data.user);
    return data.user;
  };

  const logout = () => {
    localStorage.removeItem('token');
    setUser(null);
  };

  return (
    <AuthContext.Provider value={{ user, loading, login, logout }}>
      {children}
    </AuthContext.Provider>
  );
}

export const useAuth = () => useContext(AuthContext);
\end{lstlisting}

\subsection{Intercepteur Axios}

\begin{lstlisting}[language=JavaScript,
  caption=services/api.js --- Instance Axios avec intercepteurs JWT]
import axios from 'axios';

const api = axios.create({
  baseURL: import.meta.env.VITE_API_URL
       || 'http://localhost:5000/api/v1',
});

// Injection automatique du token
api.interceptors.request.use(config => {
  const token = localStorage.getItem('token');
  if (token)
    config.headers.Authorization = 'Bearer ' + token;
  return config;
});

// Gestion globale des erreurs 401
api.interceptors.response.use(
  res  => res,
  err  => {
    if (err.response && err.response.status === 401) {
      localStorage.removeItem('token');
      window.location.href = '/login';
    }
    return Promise.reject(err);
  }
);

export default api;
\end{lstlisting}

\subsection{Routes Protégées}

\begin{lstlisting}[language=JavaScript,
  caption=components/ProtectedRoute.jsx --- Garde de route par role]
import { Navigate } from 'react-router-dom';
import { useAuth }  from '../contexts/AuthContext';

export default function ProtectedRoute({ children, roles }) {
  const { user, loading } = useAuth();

  if (loading)
    return <div>Chargement...</div>;

  if (!user)
    return <Navigate to="/login" replace />;

  if (roles && !roles.includes(user.role))
    return <Navigate to="/unauthorized" replace />;

  return children;
}
\end{lstlisting}

\begin{lstlisting}[language=JavaScript,
  caption=App.jsx --- Routes protegees par role]
<Routes>
  <Route path="/login" element={<LoginPage />} />

  <Route path="/admin/*" element={
    <ProtectedRoute roles={['admin']}>
      <AdminLayout />
    </ProtectedRoute>
  }/>

  <Route path="/teacher/*" element={
    <ProtectedRoute roles={['teacher']}>
      <TeacherLayout />
    </ProtectedRoute>
  }/>

  <Route path="/student/*" element={
    <ProtectedRoute roles={['student']}>
      <StudentLayout />
    </ProtectedRoute>
  }/>
</Routes>
\end{lstlisting}

\subsection{Page de Connexion}

\begin{lstlisting}[language=JavaScript,
  caption=pages/LoginPage.jsx --- Formulaire de connexion]
import { useState }    from 'react';
import { useNavigate } from 'react-router-dom';
import { useAuth }     from '../contexts/AuthContext';

export default function LoginPage() {
  const { login } = useAuth();
  const navigate  = useNavigate();
  const [form,  setForm]  = useState({ email: '', password: '' });
  const [error, setError] = useState('');

  const handleSubmit = async (e) => {
    e.preventDefault();
    setError('');
    try {
      const user = await login(form.email, form.password);
      if      (user.role === 'admin')   navigate('/admin');
      else if (user.role === 'teacher') navigate('/teacher');
      else                              navigate('/student');
    } catch (err) {
      const code = err.response?.data?.error;
      setError(
        code === 'INVALID_CREDENTIALS'
          ? 'Email ou mot de passe incorrect.'
          : 'Une erreur est survenue. Veuillez reessayer.'
      );
    }
  };

  return (
    <div className="min-h-screen flex items-center justify-center">
      <form onSubmit={handleSubmit}>
        <h1>Connexion -- SGUI</h1>
        {error && <p>{error}</p>}
        <input type="email"
               value={form.email}
               onChange={e =>
                 setForm({...form, email: e.target.value})} />
        <input type="password"
               value={form.password}
               onChange={e =>
                 setForm({...form, password: e.target.value})} />
        <button type="submit">Se connecter</button>
      </form>
    </div>
  );
}
\end{lstlisting}

% ────────────────────────────────────────────────────────────
\section{Sécurité Applicative}
% ────────────────────────────────────────────────────────────

\begin{table}[H]
\centering
\caption{Mesures de sécurité implémentées dans le module}
\begin{tabularx}{\linewidth}{@{}lLL@{}}
\toprule
\textbf{Menace}        & \textbf{Risque}                    & \textbf{Contre-mesure}               \\
\midrule
Injection SQL          & Acces non autorise a la BDD        & Prisma ORM (requetes parametrees)    \\
Force brute            & Compromission de compte            & Rate-limiting 20 req/15 min par IP   \\
XSS                    & Injection de scripts               & React echappe le HTML + en-tetes CSP \\
CSRF                   & Actions non autorisees             & JWT en en-tete (non envoye auto)     \\
JWT falsifie           & Usurpation d'identite              & Signature HMAC-SHA256, cle 256 bits  \\
Fuite de donnees       & Exposition du hash de mot de passe & select Prisma explicite sur chaque query \\
Elevation de role      & Acces ressources non autorisees    & requireRole sur chaque route sensible \\
\bottomrule
\end{tabularx}
\end{table}

\subsection{Configuration CORS}

\begin{lstlisting}[language=JavaScript,
  caption=app.ts --- Configuration CORS securisee]
import cors from 'cors';

app.use(cors({
  origin: process.env.FRONTEND_URL || 'http://localhost:5173',
  credentials:    true,
  methods:        ['GET','POST','PUT','PATCH','DELETE','OPTIONS'],
  allowedHeaders: ['Content-Type','Authorization'],
}));
\end{lstlisting}

% ────────────────────────────────────────────────────────────
\section{Interfaces Utilisateur}
% ────────────────────────────────────────────────────────────

\subsection{Page de Connexion}

\begin{figure}[H]
\centering
\begin{tikzpicture}
  \draw[SoftEspresso, thick, rounded corners=8pt, fill=GoldenGlow!5]
    (0,0) rectangle (9,6.5);

  % Bande superieure
  \fill[GoldenGlow] (0,5.8) rectangle (9,6.5);
  \node[font=\small\bfseries\color{SoftEspresso}] at (4.5,6.15)
    {Systeme de Gestion Universitaire — SGUI};

  % Titre
  \node[font=\large\bfseries\color{SoftEspresso}] at (4.5,5.2)
    {Connexion};
  \draw[GoldenGlow,thick] (1.5,4.9) -- (7.5,4.9);

  % Champ email
  \node[font=\scriptsize\color{AccentGray},left] at (3.5,4.6)
    {Adresse e-mail};
  \draw[AccentGray, rounded corners=4pt] (1.2,4.0) rectangle (7.8,4.4);
  \node[font=\small\color{AccentGray}] at (4.5,4.2)
    {exemple@univ-tiaret.dz};

  % Champ mot de passe
  \node[font=\scriptsize\color{AccentGray},left] at (3.5,3.7)
    {Mot de passe};
  \draw[AccentGray, rounded corners=4pt] (1.2,3.1) rectangle (7.8,3.5);
  \node[font=\small\color{AccentGray}] at (4.5,3.3)
    {\textbullet\textbullet\textbullet\textbullet\textbullet\textbullet\textbullet\textbullet};

  % Bouton
  \fill[GoldenGlow, rounded corners=4pt] (1.2,2.2) rectangle (7.8,2.7);
  \node[font=\bfseries\color{SoftEspresso}] at (4.5,2.45)
    {Se connecter};

  % Pied de page
  \node[font=\scriptsize\color{AccentGray}] at (4.5,0.4)
    {Acces securise --- SGUI v1.0 -- 2025};
\end{tikzpicture}
\caption{Maquette de la page de connexion}
\label{fig:login_ui}
\end{figure}

\subsection{Tableau de Bord Administrateur}

\begin{figure}[H]
\centering
\begin{tikzpicture}
  \draw[SoftEspresso, thick, rounded corners=6pt, fill=GoldenGlow!5]
    (0,0) rectangle (12,5.5);

  % Sidebar
  \fill[SoftEspresso] (0,0) rectangle (2.8,5.5);
  \node[font=\small\bfseries\color{GoldenGlow}] at (1.4,5.1)
    {SGUI Admin};
  \foreach \y/\lbl in {
    4.5/Tableau de bord,
    3.9/Utilisateurs,
    3.3/PFE,
    2.7/{Reclamations},
    2.1/Annonces}{
    \node[font=\scriptsize\color{white}, right] at (0.3,\y) {\lbl};
  }

  % En-tete
  \fill[GoldenGlow!20] (2.8,4.9) rectangle (12,5.5);
  \node[font=\small\bfseries\color{SoftEspresso}] at (7.4,5.2)
    {Tableau de bord --- Administrateur};

  % Cartes stats
  \foreach \x/\n/\lbl in {
    4.2/42/{Etudiants},
    7.0/18/{Enseignants},
    9.8/7/{PFE actifs}}{
    \draw[GoldenGlow, rounded corners=4pt, fill=GoldenGlow!15]
      (\x-1.1,3.2) rectangle (\x+1.1,4.6);
    \node[font=\Large\bfseries\color{SoftEspresso}] at (\x,4.1) {\n};
    \node[font=\scriptsize\color{AccentGray}]        at (\x,3.45) {\lbl};
  }

  % Tableau recent
  \draw[gray!30, fill=white, rounded corners=3pt]
    (3.0,0.3) rectangle (11.8,2.9);
  \node[font=\scriptsize\bfseries\color{SoftEspresso}] at (7.4,2.65)
    {Dernieres Reclamations};
  \draw[gray!30] (3.0,2.4) -- (11.8,2.4);
  \foreach \y/\nom/\stat in {
    2.1/{Fatima B. --- Note incorrecte}/En\ attente,
    1.7/{Karim M. --- Absence medicale}/Acceptee,
    1.3/{Sara H. --- Emploi du temps}/Rejetee}{
    \node[font=\scriptsize, right] at (3.1,\y) {\nom};
    \node[font=\scriptsize, left]  at (11.7,\y) {\stat};
  }
\end{tikzpicture}
\caption{Tableau de bord administrateur}
\label{fig:admin_ui}
\end{figure}

% ────────────────────────────────────────────────────────────
\section{Tests et Validation}
% ────────────────────────────────────────────────────────────

\subsection{Tests Fonctionnels — Authentification}

\begin{table}[H]
\centering
\caption{Résultats des tests fonctionnels — Authentification}
\begin{tabularx}{\linewidth}{@{}lLll@{}}
\toprule
\textbf{ID} & \textbf{Scenario} & \textbf{Resultat attendu} & \textbf{Statut} \\
\midrule
TF-01 & Connexion avec identifiants valides          & 200 + token JWT              & PASS \\
TF-02 & Connexion avec mot de passe incorrect        & 401 INVALID\_CREDENTIALS     & PASS \\
TF-03 & Connexion avec e-mail inexistant             & 401 INVALID\_CREDENTIALS     & PASS \\
TF-04 & Acces route protegee sans token              & 401 TOKEN\_MISSING           & PASS \\
TF-05 & Acces route protegee avec token expire       & 401 TOKEN\_INVALID\_OR\_EXPIRED & PASS \\
TF-06 & Acces route admin avec role enseignant       & 403 FORBIDDEN                & PASS \\
TF-07 & Recuperation profil \texttt{/auth/me}        & 200 + objet user             & PASS \\
TF-08 & Deconnexion (suppression token cote client)  & 200 + redirection login      & PASS \\
\bottomrule
\end{tabularx}
\end{table}

\subsection{Tests Fonctionnels — Contrôle d'Accès RBAC}

\begin{table}[H]
\centering
\caption{Résultats des tests RBAC}
\begin{tabularx}{\linewidth}{@{}lLll@{}}
\toprule
\textbf{ID} & \textbf{Scenario} & \textbf{Resultat attendu} & \textbf{Statut} \\
\midrule
TC-01 & Admin cree un compte enseignant              & 201 Created         & PASS \\
TC-02 & Enseignant tente de creer un compte          & 403 FORBIDDEN       & PASS \\
TC-03 & Etudiant soumet une reclamation              & 201 Created         & PASS \\
TC-04 & Etudiant accede a l'inbox admin              & 403 FORBIDDEN       & PASS \\
TC-05 & Enseignant modifie son sujet PFE (propose)   & 200 OK              & PASS \\
TC-06 & Enseignant modifie sujet PFE (valide)        & 403 FORBIDDEN       & PASS \\
TC-07 & Admin modifie tout sujet PFE                 & 200 OK              & PASS \\
TC-08 & Redirection auto selon role apres connexion  & Dashboard correct   & PASS \\
\bottomrule
\end{tabularx}
\end{table}

\subsection{Tests de Sécurité}

\begin{table}[H]
\centering
\caption{Tests de sécurité — Module Authentification}
\begin{tabularx}{\linewidth}{@{}lLl@{}}
\toprule
\textbf{Test} & \textbf{Description} & \textbf{Resultat} \\
\midrule
Brute-force /login &
  50 requetes en 10 secondes sur le meme compte &
  429 apres 20 requetes \\
JWT falsifie &
  Modification d'un bit dans la signature &
  401 TOKEN\_INVALID \\
Elevation de privilege &
  Token student utilise sur route admin &
  403 FORBIDDEN (100\%) \\
Injection SQL &
  Payload OR 1=1 dans le champ email &
  Neutralise (Prisma) \\
\bottomrule
\end{tabularx}
\end{table}

\subsection{Couverture de Tests}

\begin{figure}[H]
\centering
\begin{tikzpicture}
\begin{axis}[
  ybar,
  width=10cm,
  height=6cm,
  bar width=0.6cm,
  xlabel={Module},
  ylabel={Couverture (\%)},
  symbolic x coords={Auth,Users,PFE,Requests,Annonces},
  xtick=data,
  ymin=0, ymax=100,
  ymajorgrids=true,
  grid style={dashed, gray!30},
  nodes near coords,
  nodes near coords align={vertical},
  every node near coord/.append style={font=\small},
]
\addplot[fill=GoldenGlow!70, draw=SoftEspresso] coordinates {
  (Auth,     92)
  (Users,    88)
  (PFE,      85)
  (Requests, 90)
  (Annonces, 78)
};
\end{axis}
\end{tikzpicture}
\caption{Couverture de tests par module du SGUI}
\label{fig:coverage}
\end{figure}

% ────────────────────────────────────────────────────────────
\section{Conclusion}
% ────────────────────────────────────────────────────────────

Ce chapitre a présenté l'implémentation complète du module
\textbf{Authentification et Gestion des Rôles}. Le backend repose sur
un service JWT conforme au standard RFC~7519, deux middlewares chainés
(\texttt{requireAuth} + \texttt{requireRole}) qui sécurisent chaque
route de l'API, et un hachage bcrypt au coût~12 pour les mots de passe.

Le frontend utilise le Context API de React pour centraliser l'état
d'authentification, un intercepteur Axios pour l'injection transparente
du token, et des composants \texttt{ProtectedRoute} pour bloquer l'accès
aux pages non autorisées selon le rôle de l'utilisateur.

La campagne de tests fonctionnels (\textbf{16 scénarios, 100~\% PASS})
et de sécurité confirme la robustesse du système face aux principales
vulnérabilités web (injection SQL, XSS, CSRF, force brute, élévation
de privilèges). Le module constitue la fondation sécurisée sur laquelle
s'appuient tous les autres modules du SGUI.

% ============================================================
\end{document}
% ============================================================
